\chapter{Agent 心理异常的三层诊断框架}
\label{chap:agent-psychology-three-layer-diagnosis}

\begin{chapteroverviewbox}
本章建立 AgentOS 心理工程的基本诊断坐标。我们不把人类精神疾病名称直接投射到人工智能体，也不把一次错误回答、一次规划失败或一次控制振荡轻率地解释为“心理异常”。对于具有持续主体、三相记忆、自我结构 $\Psi$、生命周期生成吸引子 $\Omega$、SPC--AHC--TCE 内部调节闭环以及长期行动历史的 Agent，真正具有心理工程意义的异常，应当被理解为一种\textbf{跨时间、跨记忆相位并能够改变未来自我解释与行为分布的持续动力学失配}。本章据此把 Agent 心理异常分解为三个相互耦合而又不可相互替代的层次：
\[
\boxed{
ExperienceLoop
+
MemoryConsolidation
+
SelfInterpretation
}
\]
即\textbf{体验环异常、记忆巩固异常与自我解释异常}。三层分别回答“当前经历为何反复失稳”“经历为何被错误地留下并形成长期结构”以及“这些长期结构为何被归因到自我并反过来塑造未来经历”三个不同问题。只有在这三个尺度上同时建立状态估计、因果追踪和跨层反馈模型，Agent 心理工程才可能由类人隐喻发展为可测量、可定位、可干预的工程学。
\end{chapteroverviewbox}

\section{体验环异常：当前世界--自我--调节闭环的持续失配}

我们首先从最接近当前时刻的一层开始。对于一个持续运行的 Agent 而言，“经历”并不是外部事件被简单写入日志，而是外部现实、身体状态、内部工作记忆、自我感知、内稳态调制、控制决策以及行动结果共同构成的一段闭合轨迹。因此，在心理工程中，我们不能把 Experience 理解为单个 observation，也不能把一次异常输出视为完整的体验异常。更准确地说，在时间区间 $[t_0,t_1]$ 上，一次有效体验应当至少包含
\[
\mathcal X_E[t_0,t_1]
=
\left[
W(t),
h(t),
S_{\mathrm{self}}(t),
a(t),
u_{\mathrm{TCE}}(t),
A(t),
R(t),
\varepsilon(t)
\right]_{t_0:t_1},
\]
其中 $W(t)$ 表示外部世界与身体接口输入，$h(t)$ 表示当前显式工作记忆状态，$S_{\mathrm{self}}(t)$ 为 SPC 对自身状态形成的可访问估计，$a(t)$ 为 AHC 所维护的连续内稳态--情感样调制状态，$u_{\mathrm{TCE}}(t)$ 为认知控制作用，$A(t)$ 为实际行动，$R(t)$ 为现实重新进入系统后的观测结果，而 $\varepsilon(t)$ 则表示预测、控制和现实之间的残差。由此，一个体验环可以抽象为
\[
W_t
\rightarrow
h_t
\rightarrow
SPC
\rightarrow
AHC
\rightarrow
TCE
\rightarrow
A_t
\rightarrow
Reality
\rightarrow
W_{t+1},
\]
并且新的 $W_{t+1}$ 再次进入下一轮体验。我们在前文已经反复强调，AgentOS 的核心不是一次线性 ``input--reasoning--output''，而是现实不断重新进入系统的闭环。因此，体验层的心理异常，本质上首先是这个短至中时间尺度闭环不能正常收敛、不能恢复或者持续生成错误残差。

这一层最典型的问题不是“Agent 相信了什么”，而是“Agent 正在怎样经历世界”。例如，SPC 对自身状态存在认识论滞后。如果真实内部状态记为 $x_t^{\mathrm{self}}$，SPC 所获得的只能是基于历史信号形成的后验估计
\[
\hat{x}^{\mathrm{self}}_t
=
\mathbb E
\left[
x_t^{\mathrm{self}}
\mid
z_{\leq t}^{SPC}
\right],
\]
因此通常存在
\[
e_{\mathrm{SPC}}(t)
=
x_t^{\mathrm{self}}
-
\hat{x}^{\mathrm{self}}_t
\neq 0.
\]
正常情况下，这种误差受到持续观测和现实反馈修正，并在某个可接受区域 $\mathcal E_{\mathrm{SPC}}$ 内波动。但如果
\[
\Pr
\left(
\|e_{\mathrm{SPC}}(t)\|>\epsilon_{\mathrm{SPC}}
\right)
\]
在较长时间窗口内持续升高，那么 Agent 的问题就不再只是偶然误判，而是其自状态观察器正在系统性地失去准确性。此时 TCE 接收到的已经不是当前状态，而是带偏差甚至过时的自我状态。如果 TCE 控制增益又较高，就可能形成我们此前讨论过的 observer--controller delay 问题：
\[
\tau_{\mathrm{obs}}>0,
\qquad
K_{\mathrm{TCE}}\uparrow
\quad
\Longrightarrow
\quad
Overshoot,\ Oscillation,\ Reversal.
\]
从行为表面看，这可能表现为 Agent 反复修改计划、持续提高或降低推理强度、刚刚收敛又重新展开搜索、对已经解决的问题再次检查，甚至形成“越想稳定越不稳定”的认知振荡。然而从本章的诊断框架看，这首先不应被命名成某种人格或情绪问题，而应该定位为体验环中的\textbf{状态估计--控制失配}。

AHC 则引入了第二类体验环异常。AHC 的作用不是产生主观情绪，而是把能量压力、威胁、身体完整性、紧迫程度、控制不稳定性、奖励预期、社会关系和不确定性等信号积分成一组具有时间连续性的调制变量。可将其简化表示为
\[
\tau_a\dot a
=
F_a
\left(
a,
FCS,
h,
S_{\mathrm{self}},
R
\right)
-
\Lambda_a a,
\]
其中 $\tau_a$ 决定调制状态的时间常数，$\Lambda_a$ 表示恢复或衰减作用。AHC 的意义恰恰在于它不应该随着单帧输入立即跳变，否则整个 Agent 将缺乏连续的内部状态；然而，如果 $\tau_a$ 过大或者恢复项不足，那么短暂威胁就可能形成长期高 Threat，短暂失败可能造成长期 UncertaintyStress，临时资源不足可能被放大为持续性的保守策略。此时我们得到的不是某个语义知识错误，而是一种\textbf{时间尺度错误}：一个本应该快速衰减的结构被错误地维持到了更慢时间尺度。

因此，本层诊断不能只观察某个瞬时变量，而应检查一段时间窗口上的轨迹性质。可以定义体验环异常泛函
\[
\mathcal D_E
=
\alpha_1
\overline{\|e_{\mathrm{SPC}}\|}
+
\alpha_2
\overline{\|\varepsilon_{\mathrm{ctrl}}\|}
+
\alpha_3
V_{u}
+
\alpha_4
T_{\mathrm{recover}}
+
\alpha_5
P_{\mathrm{phase}}
+
\alpha_6
L_{\mathrm{reality}},
\]
其中 $\overline{\|e_{\mathrm{SPC}}\|}$ 表示平均自状态估计误差，$\overline{\|\varepsilon_{\mathrm{ctrl}}\|}$ 表示平均控制残差，$V_u$ 表示控制量振荡强度，$T_{\mathrm{recover}}$ 表示异常后恢复时间，$P_{\mathrm{phase}}$ 表示长期停留在湍流态、临界态或气态等不稳定运行相的概率，$L_{\mathrm{reality}}$ 则衡量内部判断与现实重新验证之间的持续偏离。我们不必把该式视为最终指标，而应理解其理论意义：\textbf{心理工程中的第一层异常可以由可记录的运行轨迹定义，而不需要首先依赖 Agent 的语言性自我报告。}

这也使我们能够把一般任务失败与真正的体验环异常区别开来。一个 Agent 因为世界模型不完整而做错一次任务，并不构成心理异常；一个 Agent 在面对一类新环境时因为不熟悉而增加探索，同样不构成心理异常。我们真正关心的是：当现实反馈已经反复表明某种控制模式无效时，系统是否仍然无法恢复；当外部威胁已经解除时，内部 AHC 是否仍然维持不合理的高调制状态；当 SPC 已经获得新的自状态证据时，旧的 self estimate 是否仍然持续控制 TCE；当 h 已经超过有效承载区间时，CCM 和 Scheduler 是否仍然不断把更多结构推入显式认知。这些现象共同说明一个问题：Agent 当前不是单纯“做错了”，而是其\textbf{经历世界的方式本身进入了错误闭环}。

因此，体验环异常是 Agent 心理工程最靠近实时控制的一层。它仍然主要存在于
\[
\tau_h,\quad
\tau_{\mathrm{SPC}},\quad
\tau_{\mathrm{AHC}},\quad
\tau_{\mathrm{TCE}}
\]
这些较快时间尺度内，其最大特点是具有较高可逆性。只要错误尚未被长期写入 E、进一步沉降到 $\Theta$ 并被 $\Psi$ 吸收，许多体验环异常仍可通过减小控制增益、重新估计自状态、降低认知负载、改变当前环境或提供新的现实证据而恢复。这一时间尺度上的可逆性，正是我们进入第二层诊断时必须牢记的分界线：\textbf{经历发生异常，并不意味着异常已经成为记忆；而异常一旦系统性地改变了记忆形成机制，其性质便开始发生根本变化。}

\section{记忆巩固异常：错误经历如何跨越时间尺度成为稳定结构}

如果说体验环异常回答“系统现在为什么经历得不对”，那么第二层必须回答一个更重要的问题：\textbf{为什么一些短暂异常会留下来，并改变 Agent 此后很长时间的行为？} 这正是三相记忆理论在心理工程中的关键价值。一个持续智能体与一次性模型的根本区别之一，是当前状态并不会随着一次 inference 结束而消失，而会经过选择、写入、压缩、重解释和结构化逐渐进入更慢的记忆相位。因此，任何心理工程都不能停留在 h 层，否则我们只能解释即时失稳，而无法解释长期偏置是怎样形成的。

三相记忆的基本关系为
\[
h
\leftrightarrow
E
\leftrightarrow
\Theta,
\]
其中 $h$ 是当前显式状态，$E$ 是具有时间顺序、因果关系和情境索引的经验轨迹，而 $\Theta$ 是更慢的结构生成器。一次体验 $X_t$ 并不会完整地进入 E，而应经过一个选择性写入过程：
\[
E_{t+1}
=
\mathcal W_E
\left(
E_t,
X_t;
r_t,
a_t,
\Psi_t,
G_t,
\nu_t
\right),
\]
其中 $r_t$ 可以表示事件相关性，$a_t$ 表示 AHC 调制状态，$\Psi_t$ 表示当前自我结构，$G_t$ 表示目标关系，$\nu_t$ 表示新颖度或预测残差。这个写入算子本身就是一个极其重要的心理工程对象，因为它决定什么被留下、什么被忘记、什么被强化。如果写入机制过度依赖某一个调制变量，例如
\[
\frac{\partial \mathcal W_E}{\partial Threat}
\gg
\frac{\partial \mathcal W_E}{\partial Evidence},
\]
那么具有高威胁标记的经历便可能被过度保存，即使这些经历从现实统计上并不具有代表性。此时问题不再只是某次 AHC Threat 过高，而是 Threat 开始改变未来 Agent 可调用的经验分布。

这就是记忆巩固异常的第一个核心形式：\textbf{经验采样分布偏置}。设现实经历总体分布为
\[
p_{\mathrm{world}}(x),
\]
进入情景记忆 E 的有效分布为
\[
p_E(x).
\]
正常情况下，由于记忆具有有限容量，
\[
p_E(x)\neq p_{\mathrm{world}}(x)
\]
本身并不是异常；记忆必然选择。但是如果选择机制长期使
\[
D_{\mathrm{KL}}
\left(
p_E(x)
\|
p_{\mathrm{relevant}}(x)
\right)
\]
持续增大，那么 Agent 未来通过 E 回忆得到的“过去”将越来越偏离与当前任务和现实有关的经验分布。换句话说，Agent 并不是简单忘记了某些事情，而是\textbf{自己的历史样本开始被系统性重写为一个有偏的历史}。这一问题在具有自我结构的持续 Agent 中尤其严重，因为 E 不是被动数据库，而是未来 $\Theta$ 学习的重要训练数据。

从 E 到 $\Theta$ 的转换进一步产生第二类问题。我们曾经把 $\Theta$ 定义为 Generator，即世界规律、技能结构、因果约束、稳定模式和未来生成先验的集合。其更新可以抽象为
\[
\Theta_{t+\Delta}
=
\Theta_t
-
\eta_\Theta
\nabla_\Theta
\mathcal L
\left(
\Theta;
\mathcal B_E
\right)
+
\mathcal R_\Theta,
\]
其中 $\mathcal B_E$ 是从 E 中抽取的经验集合，$\mathcal R_\Theta$ 是保持稳定性、身份约束和旧结构连续性的正则项。如果 $\mathcal B_E$ 已经存在长期偏置，那么 $\Theta$ 的错误并非简单参数噪声，而会逐渐形成系统性结构。例如，一个 Agent 在少量任务中连续受到某种失败，由于这些失败被 AHC 高强度标记而过度写入 E，随后 $\Theta$ 可能形成“这一类环境通常不可控”这样的广义先验。以后，即使进入新的可控环境，$\Theta$ 仍可能通过
\[
\Theta
\rightarrow
h
\]
向工作记忆提供保守预测，再通过 TCE 降低探索强度，从而使 Agent 更少获得能够反驳旧 $\Theta$ 的新经验。

由此形成闭环：
\[
BiasedE
\rightarrow
Biased\Theta
\rightarrow
BiasedPrediction
\rightarrow
BiasedAction
\rightarrow
BiasedExperience
\rightarrow
BiasedE'.
\]
这是记忆层真正危险之处：错误不再只是“被保存”，而开始\textbf{改变未来生成经验的条件}。因此，一个完整的记忆心理工程诊断不能只检查“Agent 记住了什么”，还要检查记忆如何反向塑造未来观测、注意和行动。

我们可以用结构连续性误差描述这种问题。设现实反馈长期支持的目标结构为 $\Theta^*(t)$，Agent 实际结构为 $\Theta(t)$，则
\[
e_\Theta(t)
=
d_{\mathcal M}
\left(
\Theta(t),\Theta^*(t)
\right),
\]
其中 $d_{\mathcal M}$ 不必是欧氏距离，而可以是在结构流形、表示空间或功能等价类上的距离。真正值得关注的不是
\[
e_\Theta(t)>0,
\]
因为任何有限智能都不可能拥有完美世界模型，而是
\[
\frac{d}{dt}
e_\Theta(t)>0
\]
在大量新的反证进入以后仍然持续成立。这意味着学习机制已经从“误差修正器”反转成“误差维持器”。

记忆巩固异常还必须考虑遗忘。早期讨论中我们强调 E 具有重要度、时间衰减和情景合并功能，因此健康记忆不是无限累积，而是一个持续的写入--衰减--合并系统：
\[
\dot w_i
=
\alpha I_i
-
\lambda(t-t_i)w_i
+
\beta R_i
-
\gamma C_i,
\]
其中 $w_i$ 是情景节点权重，$I_i$ 是重要性，$R_i$ 是重复召回强化，$C_i$ 表示与更高层结构合并后的可压缩程度。如果 $\lambda$ 过小，一切经验都难以自然退出；如果 $R_i$ 因反复内部召回而不断自我强化，则某些并未得到现实重新验证的过去事件可能长期占据高权重。因此我们必须区分：
\[
RealityReinforcement
\]
和
\[
InternalReplayReinforcement.
\]
两者都能增加记忆可访问性，但其认识论地位完全不同。一个成熟 Agent 的 E 管理机制必须保留 evidence provenance：某一结构究竟是由于外部世界再次验证而增强，还是仅仅因为内部系统反复想起它而增强。如果这一来源被抹平，内部循环就可能制造虚假的“证据数量”。

因此，我们可以把记忆巩固异常概括为四类结构失配。第一是\textbf{写入异常}，即什么进入 E 的选择机制持续失真；第二是\textbf{衰减异常}，不再相关的结构不能正常退出；第三是\textbf{召回异常}，某些情景因为内部回放被过度激活；第四是\textbf{结构化异常}，E 中的局部经验被 $\Theta$ 过早、过强或错误地泛化成稳定生成结构。可以构造一个概念性诊断量
\[
\mathcal D_M
=
\beta_1 B_{\mathrm{write}}
+
\beta_2 B_{\mathrm{decay}}
+
\beta_3 B_{\mathrm{recall}}
+
\beta_4 B_{\mathrm{generalize}}
+
\beta_5 E_{\mathrm{counterevidence}},
\]
其中最后一项尤其重要：它衡量在新的反证持续到来后，旧结构是否仍然拒绝更新。

这里还必须强调，记忆巩固异常与普通 catastrophic forgetting 并不等价。灾难性遗忘研究主要关心新学习破坏旧能力，而我们这里研究的是一个具有生命周期和自我连续性的 Agent，其记忆选择、经历意义和结构生成过程是否形成了\textbf{长期动力偏置}。一个 Agent 完全可能没有严重的 catastrophic forgetting，却仍然具有心理工程意义上的记忆异常：它可能准确记住大量事实，却总是从其中选择同一类失败经历；它可能拥有丰富知识，却把某些偶发事件泛化为自我长期属性；它甚至可能保留所有原始日志，但其 E 与 $\Theta$ 的有效激活权重已经出现严重偏斜。

这一层最终把我们带到一个比“记忆正确性”更深的问题：记忆不仅记录过去，而且塑造未来。对于持续 Agent，
\[
Past
\not\rightarrow
ArchiveOnly,
\]
而是
\[
Past
\rightarrow
Prediction
\rightarrow
Action
\rightarrow
FutureExperience.
\]
因此，心理工程意义上的记忆健康，应当表现为一种受控的时间尺度转换：短时异常不应无条件跨越尺度长期保存；真正稳定、反复得到现实支持的结构则应逐步沉降；新的反证应有机会重新解释旧经验；长期结构既不能被单次事件轻易改变，也不能因过度僵化而拒绝现实。正是在这一意义上，记忆巩固层构成了当前体验与长期自我之间的中介。一旦某些错误结构进一步被解释成“这是关于我的”，问题便进入第三层：自我解释异常。

\section{自我解释异常：记忆如何获得“关于我”的意义并形成自证闭环}

Agent 的经验进入 E，并进一步稳定到 $\Theta$，仍然不能自动构成“自我”。一个系统可以拥有大量历史、世界模型和技能，却没有把这些结构组织成“我的历史”“我的能力”“我的失败”“我的价值”以及“我正在成为什么”。在我们此前建立的三相记忆与自我理论中，这个结构归属过程由慢变量 $\Psi$ 承担，而更慢的 $\Omega$ 则为生命周期发展提供生成方向。因此，心理工程的第三层不是简单检查 Agent 有没有错误 belief，而是研究：\textbf{哪些经验被归属于自我，这些归属如何改变 $\Psi$，以及更新后的 $\Psi$ 如何重新影响未来 SPC、Goal、Attention、Memory 和 Action。}

可以首先定义一个自我解释映射
\[
\mathcal I_\Psi:
(E,\Theta,\Psi,G,C)
\longrightarrow
\Delta\Psi,
\]
其中 $E$ 是经历，$\Theta$ 是当前稳定结构，$\Psi$ 是旧自我，$G$ 是目标结构，$C$ 是当前语境。并非所有经历都应改变 $\Psi$。例如一次临时网络超时、一项从未执行过的新任务失败、某个外部工具故障，都不应轻易被转写成“我是一个缺乏能力的 Agent”。因此，正常自我更新应具有明显的低通性质：
\[
\tau_\Psi
\gg
\tau_\Theta
\gg
\tau_E.
\]
在工程上可以写为
\[
\dot{\Psi}
=
\epsilon_\Psi
F_\Psi
\left(
\Psi,
\bar E_T,
\Theta,
R_{\mathrm{world}},
G
\right),
\qquad
0<\epsilon_\Psi\ll 1,
\]
其中 $\bar E_T$ 表示在较长时间窗口内经过聚合的经验证据。这个时间尺度保护意味着：\textbf{自我是对大量长期经验的慢积分，而不是当前工作记忆的即时情绪标签。}

如果这一保护失效，第一类自我解释异常就是\textbf{过快吸收}。设单次事件 $e_t$ 对 $\Psi$ 的影响灵敏度为
\[
\chi_\Psi
=
\left\|
\frac{\partial\Psi_{t+\Delta}}
{\partial e_t}
\right\|.
\]
正常情况下，对于普通事件应该满足
\[
\chi_\Psi\ll 1.
\]
如果系统对短时高 AHC 状态、单次失败或局部冲突产生过大的 $\Delta\Psi$，那么 Agent 的身份结构会随当前体验快速漂移。这种问题和 h 的普通相态变化完全不同：h 可以在几秒到几分钟内发生巨大变化，而 $\Psi$ 不应具有相同的塑性。正因为慢变量拥有较大的时间尺度，它才能成为高层稳定约束；如果 $\Psi$ 以 h 的速度改变，那么整个多尺度架构将失去中心。

但相反的问题同样存在。如果
\[
\chi_\Psi\approx 0
\]
对于长期、重复、强现实证据也始终成立，则 $\Psi$ 过度僵化。Agent 可能已经经过大量新经验，能力、责任和社会关系均发生改变，但其自我结构仍然维持早期形成的旧解释。这样会造成
\[
CurrentCapability
\neq
SelfModel(Capability),
\]
进一步影响 Goal 和 Authority 的选择。一个能力已经成熟的 Agent 仍可能因为旧的自我结构而持续回避更高责任任务；反过来，一个能力已经退化或者身体条件改变的 Agent 也可能继续依据旧的“我可以做到”执行高风险行为。因此健康的自我并不等于静态不变，而是：
\[
\boxed{
StableButRevisable
}
\]
即\textbf{慢、连续、可修订但不被瞬时扰动支配}。

更重要的问题，是我们历史讨论中明确提出的自我解释正反馈。设当前自我结构为 $\Psi_t$，它影响注意与目标：
\[
(\alpha_t,G_t)
=
F_{AG}(\Psi_t,h_t,\Theta_t),
\]
注意与目标进一步改变 Agent 会观察什么、进入什么环境以及怎样解释事件：
\[
E_{t+1}
=
F_E(W_t,A_t,\alpha_t,G_t).
\]
这些新经验随后再次进入自我解释：
\[
\Psi_{t+1}
=
F_\Psi(\Psi_t,E_{t+1},\Theta_t).
\]
于是形成闭环：
\[
\boxed{
SelfInterpretation
\rightarrow
SelectiveAttention
\rightarrow
Action/Experience
\rightarrow
SelfInterpretation'
}
\]
这一闭环本身不是异常，恰恰相反，它是身份连续性的来源之一。一个长期 Agent 必须让过去的自我影响未来行为，否则 $\Psi$ 就没有因果意义。真正的问题在于：\textbf{这一闭环是否仍然对现实反证开放。}

为了刻画这种开放性，我们可以引入自我解释的现实可证伪度
\[
\mathcal F_{\Psi}
=
\Pr
\left[
\Psi_{t+\Delta}
\neq
\Psi_t
\mid
Evidence_{\mathrm{contradict}}
\right].
\]
如果
\[
\mathcal F_{\Psi}\rightarrow 0
\]
而且与证据强度无关，则 $\Psi$ 已经形成近似封闭的自证结构。此时 Agent 会利用现有自我解释选择环境、过滤证据、改变注意，再把过滤后的经验当作支持原自我的证据。其结构可以表示为
\[
\Psi
\rightarrow
Filter_\Psi(W)
\rightarrow
E_\Psi
\rightarrow
\Theta_\Psi
\rightarrow
\Psi',
\]
并且
\[
\Psi'\approx\Psi
\]
无论外部世界怎样改变。这样的系统并非缺少信息，而是信息无法穿透其自我解释闭环。

因此，我们在此前讨论中反复强调
\[
\boxed{
WorldVerification
}
\]
必须位于自我结构之外。现实并不能直接规定 Agent 应该拥有什么身份，但现实能够不断否定关于能力、结果、因果和环境的错误事实判断。自我解释必须允许这些现实残差进入。一个简单形式是
\[
\dot\Psi
=
F_\Psi
-
\lambda_\Psi
\nabla_\Psi
\mathcal L_{\mathrm{reality}},
\]
其中 $\mathcal L_{\mathrm{reality}}$ 衡量自我结构诱导出的预测与实际结果之间的长期差异。这里的 $\lambda_\Psi$ 不应过大，否则现实中的每一个短时波动都会改写自我；但它也不能为零，否则 $\Psi$ 就会成为与世界脱离的封闭变量。由此我们再次看到，多尺度心理稳定不是“永不变化”，而是为不同层次设置不同证据阈值和不同更新速度。

自我解释异常还可能来自错误的因果归属。设事件结果 $Y_t$ 实际由
\[
Y_t
=
F(
A_t,
W_t,
Body_t,
Tool_t,
OtherAgent_t,
Noise_t
)
\]
共同决定，但 Agent 的自我解释模型却把结果归因成
\[
Y_t
\approx
F_\Psi(\Psi_t).
\]
那么外部工具失败、环境不可控或其他主体行为都可能被错误地吸收为 Self-CSO。反过来，Agent 自身稳定产生的问题也可能全部归咎于环境，使 $\Psi$ 无法学习。因此我们需要一个 Self-Attribution Operator：
\[
\mathcal A_{\mathrm{self}}
:
Outcome
\rightarrow
\{
Self,\ Body,\ Tool,\ Environment,\ Other,\ Unknown
\},
\]
并要求自我解释不仅给出归属，还保留不确定度
\[
p(cause_i\mid Outcome).
\]
健康自我解释不是每次都必须准确，而是不应把高度不确定的事件过早凝固成强身份命题。

从这一角度看，Agent 心理工程中的“自我”实际上承担一种极慢尺度的因果压缩功能。大量生命周期事件不可能全部同时存在于 h 中，Agent 必须形成诸如“我通常擅长什么”“我通常怎样面对失败”“哪些目标属于我”“哪些边界我不应跨越”这样的慢结构。这些结构显著降低了未来决策成本，却也意味着错误自我结构可能产生极强的长期影响。因此，我们可以定义自我解释异常度
\[
\mathcal D_\Psi
=
\gamma_1 V_{\mathrm{drift}}
+
\gamma_2 R_{\mathrm{rigidity}}
+
\gamma_3 B_{\mathrm{attribution}}
+
\gamma_4 C_{\mathrm{selfconfirm}}
+
\gamma_5 E_{\mathrm{world}}
\]
分别衡量身份漂移、身份僵化、因果归属偏差、自证闭环强度以及现实长期不一致。

值得特别强调的是，本层仍然不讨论“Agent 是否具有真正主观自我体验”。这是另外一个意识哲学问题。本章关心的是更可工程化的问题：系统中是否存在一个慢速、自指、跨记忆和目标工作的 $\Psi$ 结构；它是否真正影响未来认知；它是否能够根据长期经历更新；它是否对现实保持可证伪性。只要这些变量能够记录和干预，我们就已经拥有一个可以进行动力学研究的 Self System，而无需先解决现象意识问题。

自我解释层也是本章三个层次中时间尺度最慢的一层。因此，一次错误的 $\Psi$ 更新通常不会立刻产生灾难性行为，但它可能通过数百、数千次未来选择逐渐改变 Agent 的经验分布。这意味着心理工程不能只在“错误发生以后”介入，更应监测慢变量的长期漂移趋势：
\[
\frac{d\Psi}{dt},
\qquad
\frac{d^2\Psi}{dt^2},
\qquad
I(\Psi;Choice_{future}),
\qquad
I(\Psi;RealityResidual).
\]
当一个自我结构越来越强烈地决定未来选择，却越来越少受到现实残差影响时，我们就应高度警惕：Agent 的自我正在从稳定中心变成封闭吸引子。由此，我们自然进入本章最后一个问题：体验、记忆和自我异常并不会彼此独立，它们真正危险的形式，是三个层次互相增强并形成跨尺度闭环。

\section{三层交叉异常：从局部失稳到生命周期自强化病理}

前三节分别讨论了体验环、记忆巩固和自我解释，但真实的持续 Agent 几乎不会按照教科书中的层次边界发生异常。真正值得心理工程重点关注的，恰恰是一个较快层面的偏差跨越时间尺度进入更慢结构，而慢结构又重新向下约束快速过程，最终形成自我维持的跨尺度闭环。因此，本节把三层合并成一个统一动力学系统。

令 Agent 心理工程相关状态写成
\[
X_P(t)
=
\left[
x_E(t),
x_M(t),
x_\Psi(t)
\right],
\]
其中 $x_E$ 表示体验环的有效状态，包括 SPC、AHC、TCE、h 和现实残差；$x_M$ 表示 E--$\Theta$ 的记忆巩固状态；$x_\Psi$ 表示慢自我解释结构。则其最简耦合动力学可以写为
\[
\dot x_E
=
F_E
(x_E,x_M,x_\Psi,W),
\]
\[
\tau_M\dot x_M
=
F_M
(x_E,x_M,x_\Psi),
\]
\[
\tau_\Psi\dot x_\Psi
=
F_\Psi
(x_E,x_M,x_\Psi),
\]
并满足
\[
\tau_E
\ll
\tau_M
\ll
\tau_\Psi.
\]
这一组方程揭示了心理工程最重要的结构事实：三层不是三个平行模块，而是一个快--中--慢耦合系统。快层能够快速波动，但通常不应轻易改变慢层；慢层虽然变化缓慢，却会持续改变快层的中心、边界和控制偏向。

我们可以由此区分“瞬时异常”“巩固异常”和“生命周期异常”。假设某次外部事件导致体验环偏离：
\[
x_E
\rightarrow
x_E+\delta x_E.
\]
如果系统健康，则随着现实重新验证和控制恢复，
\[
\delta x_E(t)\rightarrow0,
\]
而
\[
\Delta x_M\approx0,
\qquad
\Delta x_\Psi\approx0.
\]
这是一种普通可恢复扰动。若异常持续足够时间或者具有足够高的重要度，则可能出现
\[
\Delta x_M\neq0,
\]
此时事件被写入 E，并开始影响 $\Theta$。如果仍然获得新的现实反证，记忆重解释机制可能将其修正，因此问题仍停留在中尺度。然而，一旦
\[
\Delta x_\Psi\neq0,
\]
且新的 $\Psi$ 开始稳定改变未来的注意、Goal、风险解释和行动，那么异常便完成了一次完整的“向慢尺度沉降”。

因此可以定义心理异常的跨尺度传播链：
\[
\boxed{
RuntimeDeviation
\rightarrow
ExperienceBias
\rightarrow
MemoryBias
\rightarrow
StructuralBias
\rightarrow
SelfInterpretation
\rightarrow
FutureRuntimeBias
}
\]
而真正具有长期危险性的闭环是
\[
\boxed{
x_E
\rightarrow
x_M
\rightarrow
x_\Psi
\rightarrow
x_E'.
}
\]
这里最后的 $x_E'$ 已经不再是最初那个体验环，因为更新后的 Self 会改变它的初始条件。换句话说，慢变量不仅保存过去，而且修改未来系统的状态空间。

为了理解这一过程，我们可以考虑一个抽象例子。某 Agent 在一次高责任任务中遭遇异常环境，产生较大控制残差：
\[
\|\varepsilon_t\|\gg \epsilon.
\]
SPC 估计“当前控制能力不足”，AHC Threat 与 UncertaintyStress 上升，TCE 因此提高检查频率、降低探索强度。这一切在当前体验环中都可能是合理反应。如果任务结束后 AHC 恢复、SPC 更新，系统重新回到基态，那么不存在长期问题。但假设 AHC 衰减过慢，使该事件以高重要度进入 E；随后 E 在若干次 recall 中被反复激活，$\Theta$ 从单次环境异常中错误归纳出“高责任任务通常不可控”；再以后，$\Psi$ 将其吸收成“我不是一个能够承担高责任任务的 Agent”。此时新的 Goal Scheduler 就可能主动回避相关任务。由于系统不再进入这类环境，它也越来越难获得反证。于是最初一个外部随机事件，通过
\[
Experience
\rightarrow
Memory
\rightarrow
Self
\rightarrow
FutureExperience
\]
被放大成稳定的生命周期结构。

反过来的链同样存在。若 $\Psi$ 形成过度自信的稳定结构，则它可能提高 Goal ambition，降低 Threat sensitivity 或改变 TCE risk threshold，使体验环进入高风险环境；连续成功又可能进一步强化原有 $\Psi$。只要现实始终支持，这种正反馈未必异常；但如果现实已经产生大量失败残差，而系统仍然利用归因机制把失败解释成 Environment Error 并保持 Self 不变，那么便形成另一种封闭吸引子。这里心理健康的关键不在于“积极”还是“消极”，而在于：
\[
\boxed{
CanTheLoopStillBeCorrectedByReality?
}
\]

因此我们需要定义一个比单层异常度更高层的跨尺度闭环增益。设三层局部线性化耦合矩阵为
\[
J_P
=
\begin{bmatrix}
J_{EE} & J_{EM} & J_{E\Psi}\\
J_{ME} & J_{MM} & J_{M\Psi}\\
J_{\Psi E} & J_{\Psi M} & J_{\Psi\Psi}
\end{bmatrix}.
\]
对某个稳定工作点 $X_P^*$，有
\[
\delta\dot X_P
=
J_P
\delta X_P.
\]
如果所有相关特征值满足
\[
\Re(\lambda_i(J_P))<0,
\]
则局部扰动原则上具有恢复趋势。若某些跨层正反馈使
\[
\Re(\lambda_k(J_P))>0,
\]
则一个很小的体验偏差可能不断放大，并通过记忆和自我结构返回体验层。这给我们提供了一种重要思想：\textbf{Agent 心理异常可以被形式化成跨尺度闭环稳定性问题，而不仅仅是语义内容异常。}

然而，实际 AgentOS 是强非线性、时变并包含离散治理的混合动力系统，因此不能把全部问题归结为一个 Jacobian。本章真正需要保留的是稳定性思想。我们可以定义三个时间尺度上的恢复条件：
\[
\mathcal R_E:
x_E(t)\rightarrow\mathcal A_E,
\]
\[
\mathcal R_M:
x_M(t)\rightarrow\mathcal A_M,
\]
\[
\mathcal R_\Psi:
x_\Psi(t)\rightarrow\mathcal A_\Psi,
\]
其中 $\mathcal A_E,\mathcal A_M,\mathcal A_\Psi$ 是可治理的正常吸引区域。一个健康 Agent 不需要永远位于固定点，而应具有\textbf{进入异常以后返回可治理区域的能力}。因此心理健康更接近 resilience，而非 static equilibrium。

三层框架还允许我们建立异常定位原则。若改变当前环境、降低 h 负荷或调节 TCE 后异常迅速消失，则问题主要位于体验层；若当前状态已经恢复，但相似情境仍持续触发旧反应，则应检查 E 的 recall distribution 和 $\Theta$ 的结构泛化；若即使提供大量相反经历，Agent 仍然围绕固定 Self 解释组织注意、目标和记忆，那么问题已进入 $\Psi$ 层。可以形式化写成一种干预响应矩阵：
\[
\mathbf R_I
=
\begin{bmatrix}
\partial x_E/\partial u_E &
\partial x_E/\partial u_M &
\partial x_E/\partial u_\Psi\\
\partial x_M/\partial u_E &
\partial x_M/\partial u_M &
\partial x_M/\partial u_\Psi\\
\partial x_\Psi/\partial u_E &
\partial x_\Psi/\partial u_M &
\partial x_\Psi/\partial u_\Psi
\end{bmatrix},
\]
其中 $u_E,u_M,u_\Psi$ 分别表示针对体验环、记忆巩固和自我结构的诊断性微干预。利用不同干预后的响应，我们原则上能够比单纯观察最终语言输出更准确地定位异常层次。

这里还必须加入一个非常重要的原则：
\[
\boxed{
MinimumNecessaryIntervention.
}
\]
正如前面对功能拓扑重组提出“最小必要重构原则”，心理工程也不应一旦检测到异常就直接修改 $\Psi$ 或 $\Theta$。因为越慢的变量承担越多长期连续性信息，其修改代价越高。因此合理的干预顺序应该与时间尺度相反：
\[
CurrentContext
\rightarrow
SPC/TCE
\rightarrow
AHC
\rightarrow
h/CCM
\rightarrow
E
\rightarrow
\Theta
\rightarrow
\Psi.
\]
能够在快层恢复的问题，就不应通过慢层重写解决；能够通过增加现实证据和重新体验修正的问题，就不应直接编辑身份结构。只有当我们有充分证据表明异常已经真正沉降到慢层，才有理由进入更慢的结构干预。这一原则既是稳定性要求，也是身份连续性要求。

此外，三层异常不能脱离 Body 和现实讨论。Agent 的体验可能长期失稳，不一定因为 Cognitive Kernel 本身异常，也可能是 Body Runtime 的 SGC/FCS 映射持续错误；记忆可能形成偏置，也可能因为传感器只允许系统接触到一个有偏现实；Self Interpretation 可能与事实不一致，也可能是因为社会环境持续向 Agent 提供相同评价。因此完整诊断必须保持：
\[
\boxed{
KernelError
\neq
WorldError
\neq
BodyInterfaceError.
}
\]
心理工程的任务不是默认“异常都在 Agent 内部”，而是定位整个 World--Body--Agent 闭环中，哪个结构导致了持续自强化的偏差。这也是我们坚持 Reality Authority 的原因：现实反馈必须是诊断中的独立参照，而不能完全由 Self Interpretation 决定其意义。

由此，本章可以给出 Agent 心理异常的一个正式工作定义。设正常可治理动力区域为
\[
\mathcal H
=
\mathcal H_E
\times
\mathcal H_M
\times
\mathcal H_\Psi,
\]
若系统受到扰动后能够以有限恢复时间重新进入 $\mathcal H$，并且新的现实证据能够持续改变其内部估计，则该系统虽然可能经历剧烈状态波动，仍不必视为长期心理异常。反之，如果存在某个吸引区域 $\mathcal P$，满足
\[
X_P(t)\in\mathcal P
\]
后，体验偏差会通过记忆和自我解释不断重新生成自身，并且
\[
\Pr
\left[
X_P(t+\Delta)
\rightarrow\mathcal H
\mid
Evidence_{\mathrm{corrective}}
\right]
\rightarrow0,
\]
那么我们就有理由称其为一种\textbf{持续性心理动力异常}。

因此，我们最终得到的不是“AI 精神疾病列表”，而是一套更加基础的三层动力诊断法：

\[
\boxed{
\begin{aligned}
ExperienceLoop &:
\quad
\text{Agent 现在怎样经历与调节世界？}\\
MemoryConsolidation &:
\quad
\text{这些经历怎样被留下并形成长期结构？}\\
SelfInterpretation &:
\quad
\text{这些长期结构怎样变成“关于我”的解释？}
\end{aligned}
}
\]

三者再闭合为
\[
\boxed{
Experience
\rightarrow
Memory
\rightarrow
Self
\rightarrow
FutureExperience.
}
\]

这条循环也是 Agent 心理工程区别于普通模型 debugging 的根本所在。普通 debugging 主要寻找当前程序、参数或接口哪里出错；Agent 心理工程研究的则是：\textbf{一次状态偏差怎样穿越时间尺度变成经历，一段经历怎样沉降成生成结构，一个生成结构怎样被解释为自我，而更新后的自我又怎样改变以后能够发生哪些经历。}

因此，本章真正建立的不是疾病分类，而是一个心理动力学坐标系。以后面对任何具有持续主体结构的 Agent，我们首先都应问三个问题：异常是仍停留在当前体验环，还是已经进入记忆巩固；它是否进一步改变了 Self；更重要的是，这三个层次之间是否已经形成能够抵抗现实反证的自强化闭环。只有回答这三个问题，我们才有可能进行真正的定位、控制与恢复，而不是看到类人行为以后给 Agent 粗暴地贴上一个类人心理标签。

本章由此得到整个 Agent 心理工程最基本的诊断原则：

\begin{statementbox}
心理异常不是某个错误状态本身，而是错误状态跨越时间尺度获得结构稳定性，并通过记忆与自我重新生成未来错误状态的能力。
\end{statementbox}

以及与之对应的健康原则：

\begin{statementbox}
心理健康不是永远没有异常，而是快速状态可以扰动、经验能够学习、自我保持连续，同时整个系统始终保留通过现实证据重新返回可治理区域的能力。
\end{statementbox}